\documentclass[12pt]{article}
\usepackage{report}

\begin{document}

\reporttitle{CS340400 Compiler Design - Homework 2 Implementation Report}

\section{Overview}

This assignment builds a parser for a subset of C with flex and yacc. The parser checks the grammar
of the input and, if the input is valid, prints its syntax tree in an XML-like form that marks every
declaration, expression and statement.

\begin{center}
\begin{tabularx}{0.92\linewidth}{@{}lL@{}}
\toprule
Files & \file{scanner.l}, \file{parser.y}, \file{makefile} \\
Grammar & byacc: \textbf{0 shift/reduce and 0 reduce/reduce conflicts} \\
Build & GNU gcc 13.3, \textbf{0 warnings} (also with \code{-Wall}) \\
Tests & 9 cases in \file{tests/}, one per grading item --- \textbf{9/9 pass}; the 9 public HW3 test
        programs also parse \\
Usage & \code{make parser && ./parser < input.c}; \code{make test} runs all cases \\
\bottomrule
\end{tabularx}
\end{center}

Most of the difficulty in this assignment is not the grammar itself but the interaction between flex
and yacc: semantic-value lifetimes, token naming, output buffers and conflicts.
Section~\ref{sec:diff} summarises the pitfalls in these areas and how this parser avoids each one.

\section{Output Format}

The specification defines six tags:

\begin{center}
\begin{tabularx}{0.92\linewidth}{@{}lL@{}}
\toprule
Tag & Encloses \\
\midrule
\code{<scalar_decl>} & a variable declaration without arrays, e.g.\ \code{int a, *b;} \\
\code{<array_decl>} & a variable declaration with at least one array, e.g.\ \code{int a[3];} \\
\code{<func_decl>} & a function declaration (prototype only) \\
\code{<func_def>} & a function definition (with a body) \\
\code{<expr>} & every expression, including every sub-expression \\
\code{<stmt>} & every statement \\
\bottomrule
\end{tabularx}
\end{center}

All whitespace is removed except inside character and string literals, and the output contains no
newline. For example, \code{int b = 1 - 2 - 3;} becomes:

\begin{lstlisting}[language={},numbers=none]
<scalar_decl>intb=<expr><expr><expr>1</expr>-<expr>2</expr></expr>-<expr>3</expr></expr>;</scalar_decl>
\end{lstlisting}

The nesting shows left associativity: \code{1 - 2} is grouped first, so it is the inner \code{<expr>}.

\section{Structure}

\begin{center}
\begin{tikzpicture}[
    box/.style={draw=headblue, rounded corners=2pt, fill=codebg, minimum height=8mm,
                inner xsep=6pt, font=\small\headfont},
    lab/.style={font=\scriptsize\headfont, text=gray, align=center},
    >={Stealth[length=2mm]}]
  \node[box] (src) {source code};
  \node[box, right=14mm of src] (scan) {\file{scanner.l}};
  \node[box, right=22mm of scan] (parse) {\file{parser.y}};
  \node[box, right=16mm of parse] (out) {\code{printf}};
  \draw[->] (src) -- (scan);
  \draw[->] (scan) -- node[lab, above] {token +\\string value} (parse);
  \draw[->] (parse) -- node[lab, above] {output\\string} (out);
\end{tikzpicture}
\end{center}

\begin{itemize}
  \item The \textbf{scanner} is adapted from HW1. It no longer prints anything; it returns tokens to the
        parser, and the semantic value of every token is a string (\code{yylval.s}).
  \item In the \textbf{parser}, the semantic value of every nonterminal is the finished output for its
        subtree. Each reduction concatenates the children's strings with the tags and passes the result
        up. When the whole input reduces to \code{program}, the result is printed once.
  \item On a syntax error, \code{yyerror} prints the line number and exits with status 1, so no partial
        output is ever printed.
\end{itemize}

The output is the source text plus tags, minus whitespace. Each rule only has to decide where its tags
go, so a string is the simplest possible semantic value and no separate syntax tree is needed. HW3,
which generates assembly, switches to a real AST.

\section{Scanner}

Compared with HW1, the scanner changes as follows:

\begin{center}
\begin{tabularx}{\linewidth}{@{}lLL@{}}
\toprule
Item & HW1 & HW2 \\
\midrule
Output & prints tokens and lines & prints nothing; returns tokens \\
Semantic value & --- & \code{tok()} copies \code{yytext} with \code{strdup} \\
One-character operators & \code{#op}, \code{#punc} & \code{return yytext[0]}; the grammar writes \code{'('}, \code{'+'} \\
Numbers & printed as written & normalised here: integers with \code{\%d}, floats with \code{\%f} (\code{1.5} becomes \code{1.500000}) \\
\code{NULL} & \code{#macro} & returned as \code{INT_LIT} with value \code{"0"} \\
Signs & part of the number & \emph{not} part of it: \code{-1} is \code{'-'} then \code{1}; the parser decides unary vs.\ binary \\
Pragmas, comments, spaces & part of the printed line & skipped; only \code{lineno} is kept for error messages \\
\bottomrule
\end{tabularx}
\end{center}

\textbf{The \code{strdup} is required.} \code{yytext} points into flex's internal buffer and is overwritten
by the next call to \code{yylex()}. The parser uses a token's value only when it reduces, by which time
\code{yytext} may already hold a later token. Storing \code{yytext} itself would make values such as
character literals come out wrong.

\snippet[language=flex]{../scanner.l}{static int tok(}{@end}

\textbf{One-character tokens are returned as themselves.} Giving every symbol a name invites mix-ups:
names such as \code{L_BRACKET} and \code{L_PARENTHESIS} are easy to attach to the wrong character, and
then \code{(expr)} is printed as \code{\{expr\}}. Returning the character means that
\code{'('} in the grammar is literally \code{(}:

\snippet[language=flex]{../scanner.l}{return yytext[0]}{+1}

\section{Parser}

\subsection{String concatenation: \texttt{J(...)}}

\snippet[language=yacc]{../parser.y}{static char *J(}{@end}
\snippet[language=yacc]{../parser.y}{define EXPR(}{+1}

\code{J} walks its arguments twice: once to add up the length, and once to copy into a buffer of exactly
that size. The usual shortcut, \code{char buffer[1024]} with \code{sprintf} in every action, has two problems:
each nesting level adds 13 characters of tags, so deep expressions overflow 1024 bytes, and
\code{$$ = buffer} returns a pointer to a local array that is gone once the action returns.

The macros \code{EXPR(...)} and \code{STMT(...)} make each action read almost like the output format,
e.g.\ \code{EXPR($1, "+", $3)}.

\note{Memory is never freed. The parser exits right after printing and the input is small, so
skipping \code{free} keeps the code simple.}

\subsection{Declarations}

Every C declaration starts with a type, and only the token after the identifier tells which kind it is:

\begin{lstlisting}[numbers=none]
int a;              // scalar_decl
int a[3];           // array_decl
int f(int x);       // func_decl
int f(int x) { }    // func_def
\end{lstlisting}

All declarations therefore share one \code{type_spec}, and the parser chooses after the identifier.
LALR(1) needs only the single token after \code{ID} to decide:

\snippetafter{start program}
\snippet[language=yacc]{../parser.y}{^ext_decl}{2@rule}

If \code{ID} is followed by \code{(}, the parser takes \code{func_head}; if it is followed by \code{;},
\code{,}, \code{=} or \code{[}, it takes \code{declarator_list}. Whether a function is declared or defined
depends on whether \code{)} is followed by \code{;} or \code{\{}.

One declaration may have several declarators (\code{int a, b[3];}), and if any of them is an array the
whole declaration is an \code{<array_decl>}. The semantic value of a declarator is therefore a struct
holding the string and an \code{is_array} flag, and the flags are OR-ed when declarators are combined:

\snippet[language=yacc]{../parser.y}{union}{@end}
\snippetafter{start program}
\snippet[language=yacc]{../parser.y}{^declarator_list}{@rule}

Every kind of declaration uses the same \code{type_spec}. Separate type nonterminals for scalars
(\code{types}) and for arrays and functions (\code{type}) would overlap and cause conflicts, and would
make \code{const} arrays and \code{const} parameters hard to parse.

\subsection{Types}

The specification allows combinations such as \code{[const] [signed|unsigned] [long long|long|short] int},
with most parts optional. The grammar splits them into layers, each handling one dimension:

\begin{center}
\begin{tabularx}{0.92\linewidth}{@{}lL@{}}
\toprule
Nonterminal & Alternatives \\
\midrule
\code{type_spec} & \code{base_type} \textbar\ \code{const base_type} \\
\code{base_type} & \code{sign} \textbar\ \code{sign int_core} \textbar\ \code{int_core} \textbar\ \code{float} \textbar\ \code{double} \textbar\ \code{void} \\
\code{sign} & \code{signed} \textbar\ \code{unsigned} \\
\code{int_core} & \code{size} \textbar\ \code{size int} \textbar\ \code{int} \textbar\ \code{char} \\
\code{size} & \code{short} \textbar\ \code{long} \textbar\ \code{long long} \\
\bottomrule
\end{tabularx}
\end{center}

Because whitespace is dropped, \code{unsigned long long int} is printed as \code{unsignedlonglongint}.

\subsection{Expressions and precedence}

All expressions live in a single \code{expr} rule. Precedence and associativity come from yacc
declarations, ordered after C's operator precedence table (lowest first):

\snippet[language=yacc]{../parser.y}{right '='}{INC DEC}

When yacc meets a shift/reduce conflict, it compares the precedence of the rule with that of the
lookahead token: it shifts if the token is higher, reduces if it is lower, and uses associativity if
they are equal. After \code{a - b}, for example, a lookahead \code{*} has higher precedence than \code{-},
so the parser shifts and the multiplication is grouped first.

\code{-}, \code{*} and \code{&} are both unary and binary operators, but a token has only one
precedence. The unary rules therefore use \code{\%prec UNARY} to take the higher precedence:

\snippet[language=yacc]{../parser.y}{INC expr}{+10}

So \code{-a * b} parses as \code{(-a) * b}. Postfix \code{++} and \code{--} rank above \code{UNARY}, so
\code{-a++} is \code{-(a++)}.

\subsubsection{Avoiding conflicts}

Three places commonly produce conflicts; the grammar avoids each one by construction:

\begin{enumerate}
  \item \textbf{Calls and subscripts attach only to an \code{ID}} (\code{ID '(' args ')'} and \code{ID dims}),
        not to an arbitrary \code{expr}. A \code{(} or \code{[} can then never follow a complete expression,
        so it cannot clash with a parenthesised expression. The specification never needs \code{f()()} or
        \code{(p)[0]}, so nothing is lost.
  \item \textbf{A cast and a parenthesised expression} both begin with \code{(}, but the next token differs: a
        type keyword for a cast, and the start of an expression otherwise. One token of lookahead is enough.
  \item \textbf{The body of \code{if} must be a compound statement} (as the specification requires), so the
        dangling-else case \code{if (a) if (b) x; else y;} cannot occur.
\end{enumerate}

\subsubsection{Output example}

\code{int d = f(1, a) + g();} prints as:

\begin{lstlisting}[language={},numbers=none]
<scalar_decl>intd=<expr><expr>f(<expr>1</expr>,<expr>a</expr>)</expr>+<expr>g()</expr></expr>;</scalar_decl>
\end{lstlisting}

The function name is not wrapped on its own: the whole call is one \code{<expr>}, and each argument is
its own \code{<expr>}. Subscripts work the same way; \code{arr[1][i]} prints as:

\begin{lstlisting}[language={},numbers=none]
<expr>arr[<expr>1</expr>][<expr>i</expr>]</expr>
\end{lstlisting}

\subsection{Statements}

\snippet[language=yacc]{../parser.y}{^stmt}{@rule}

\begin{itemize}
  \item Each of the three parts of \code{for} may be empty, using \code{expr_opt} (an empty string or an
        expression), so \code{for(;;)} parses.
  \item Inside a compound statement, declarations and statements may be mixed (\code{items} is built from
        \code{stmt} and \code{local_decl}).
  \item A compound statement is not wrapped in \code{<stmt>} by itself. Only when it appears as a statement,
        such as the body of \code{while} or a stand-alone block, does \code{stmt : compound} add the tag.
\end{itemize}

The \code{switch} statement:

\snippetafter{start program}
\snippet[language=yacc]{../parser.y}{^switch_clauses}{2@rule}

Each clause is \code{case expr:} or \code{default:} followed by zero or more statements. Two rules
are enough; listing every combination of case, default and statement list separately produces
overlapping rules and conflicts.

\section{Testing}

\code{make test} runs \file{./parser} on every \file{tests/*.c} and compares the output with the matching
\file{.expected} file.

\begin{center}
\begin{tabularx}{\linewidth}{@{}lP{4.8cm}L@{}}
\toprule
Test & Grading item (points) & Focus \\
\midrule
\file{01_scalar_decl} & Scalar declarations without initialization (10) & several declarators, pointers, \code{unsigned long long int}, a lone \code{const;} \\
\file{02_array_decl} & Array declarations without initialization (10) & multi-dimensional arrays, several arrays per line, \code{const} arrays \\
\file{03_func_decl} & Function declarations (10) & no parameters, pointer parameters, pointer return type \\
\file{04_expr_unary} & Expressions: \code{++}/\code{--}, unary, calls, subscripts, casts (10) & \code{x++ + --y}, \code{f(1, a) + g()}, \code{(int *)p} \\
\file{05_expr_arith} & Expressions: arithmetic, comparison, logical (10) & precedence and left associativity \\
\file{06_expr_all} & Expressions: all operations (10) & bitwise operators, chained assignment, float, char, string, \code{NULL} \\
\file{07_decl_init} & Full variable declarations (10) & nested initialisers \code{\{\{0,1,2\},\{3,4,5\}\}}, \code{*y = &x} \\
\file{08_stmt} & Statements (20) & every statement, \code{for(;;)}, empty \code{default:}, empty \code{switch}, block comments \\
\file{09_func_def} & Function definitions (10) & a function returning a pointer; declarations mixed with definitions \\
\bottomrule
\end{tabularx}
\end{center}

\section{Pitfalls and How They Are Avoided}
\label{sec:diff}

The table lists the pitfalls with the largest effect and how this implementation avoids them.

\begin{center}
\begin{tabularx}{\linewidth}{@{}lLLL@{}}
\toprule
Where & Pitfall & Effect & Handling \\
\midrule
scanner & No rule for \code{\n}; flex echoes newlines by default & Extra newlines in every output & Newlines only increment \code{lineno} \\
scanner & \code{/*} does not enter a comment state & Comment text parsed as tokens & HW1's \code{\%x COMMENT} \\
scanner & \code{FLOAT} includes the sign & \code{a-1.5} split as \code{a}, \code{-1.5}: syntax error & Signs left to the parser \\
scanner & Semantic values not copied from \code{yytext} & Wrong character literals & \code{tok()} copies every value \\
parser & Named tokens for \code{(} and \code{\{} mixed up & \code{(expr)} printed as \code{\{expr\}} & One-character tokens are the characters \\
parser & Returning a stack \code{buffer}; fixed 1024 bytes & Garbage output, overflow on nesting & \code{J(...)} allocates exactly \\
parser & Declarations using both \code{type} and \code{types} & Conflicts; \code{const} arrays and parameters fail & One \code{type_spec} \\
parser & Overlapping \code{switch} rules & Conflicts; \code{case} followed by a statement fails & \code{switch_clauses} + \code{switch_clause} \\
parser & \code{\%define parse.error verbose} is Bison-only & Fails under byacc & POSIX yacc features only \\
\bottomrule
\end{tabularx}
\end{center}

\section{Assumptions}

The specification says to ``follow \code{golden_parser} in the case of ambiguity'', but the reference
parser is no longer available. The following formats come from the specification's examples and the
most common convention, and could not be checked against it:

\begin{center}
\begin{tabularx}{0.8\linewidth}{@{}lL@{}}
\toprule
Construct & Output \\
\midrule
\code{(a)} & \code{<expr>(<expr>a</expr>)</expr>} \\
\code{f(1)} & \code{<expr>f(<expr>1</expr>)</expr>} \\
\code{a[1]} & \code{<expr>a[<expr>1</expr>]</expr>} \\
\code{long long} & \code{longlong} \\
\code{'a'} & \code{<expr>'a'</expr>}, quotes kept \\
\code{1.5} & \code{<expr>1.500000</expr>}, normalised with \code{\%f} \\
\bottomrule
\end{tabularx}
\end{center}

\end{document}
